\documentclass{article}
\usepackage{enumitem}
\usepackage{amsmath}
\begin{document}

\title{Chapter 35: Organisms and Populations}
\date{}
\maketitle

\begin{enumerate}[label=Q\arabic*.]

\item The logistic growth model is represented by the equation:
\\(A) $dN/dt = rN$
\\(B) $dN/dt = rN\left(\frac{K-N}{K}\right)$
\\(C) $N_t = N_0 e^{rt}$
\\(D) $dN/dt = rN + K$

\item The carrying capacity (K) in logistic growth represents:
\\(A) The maximum growth rate of the population
\\(B) The maximum population size the environment can sustainably support
\\(C) The initial population size
\\(D) The rate of population increase

\item Which of the following is a density-dependent factor regulating population growth?
\\(A) Earthquake
\\(B) Flood
\\(C) Intraspecific competition for food
\\(D) Temperature extremes

\item The competitive exclusion principle (Gause's principle) states:
\\(A) Two species can coexist indefinitely if they have similar niches
\\(B) Two species competing for the identical ecological niche cannot coexist; one will outcompete the other
\\(C) Competition always leads to extinction of both species
\\(D) Competition promotes speciation

\item Assertion: Mutualism is a symbiotic relationship where both species benefit.\\
Reason: In mycorrhizal association, the fungus benefits from plant photosynthate and the plant benefits from enhanced mineral absorption.
\\(A) Both true, Reason is correct explanation
\\(B) Both true, Reason is not correct explanation
\\(C) Assertion true, Reason false
\\(D) Assertion false, Reason true

\item Age pyramids in human population studies where the base (young age group) is broad indicate:
\\(A) Declining population
\\(B) Stable population
\\(C) Rapidly expanding (growing) population
\\(D) Post-reproductive bulge

\item Natality refers to:
\\(A) Death rate in a population
\\(B) Birth rate; number of new individuals added per unit time
\\(C) Immigration into a population
\\(D) Age structure of the population

\item Commensalism is an interaction where:
\\(A) Both species benefit
\\(B) One species benefits and the other is unaffected
\\(C) One species benefits and the other is harmed
\\(D) Both species are harmed

\item The predator-prey oscillations (Lotka-Volterra cycles) show that:
\\(A) Predator and prey populations are completely independent
\\(B) Prey population peaks are followed by predator population peaks with a time lag
\\(C) Predators drive prey to extinction over time
\\(D) Prey populations remain stable regardless of predator numbers

\item Camouflage in animals like stick insects is an adaptation for:
\\(A) Attracting mates
\\(B) Avoiding detection by predators (crypsis)
\\(C) Thermoregulation
\\(D) Communication between individuals

\item Batesian mimicry involves:
\\(A) Two unpalatable species resembling each other (Müllerian)
\\(B) A palatable/harmless species resembling an unpalatable/dangerous species
\\(C) Both species having identical coloration for intraspecific communication
\\(D) Aggressive mimicry for predation

\item Which of the following describes the r-selected species strategy?
\\(A) Large body size, long lifespan, few offspring, high parental care
\\(B) Small body size, short lifespan, many offspring, little parental care, early reproduction
\\(C) K closer to carrying capacity, stable environments
\\(D) Population size close to carrying capacity always

\item Ectoparasites differ from endoparasites in that:
\\(A) Ectoparasites live inside the host's body
\\(B) Ectoparasites live on the external surface of the host
\\(C) Ectoparasites are more harmful than endoparasites
\\(D) Ectoparasites do not affect host fitness

\item Match the following population interactions with examples:\\
1. Mutualism $\rightarrow$ (a) Lichens (alga + fungus)\\
2. Commensalism $\rightarrow$ (b) Barnacles on whale\\
3. Parasitism $\rightarrow$ (c) \textit{Plasmodium} in human RBCs\\
4. Amensalism $\rightarrow$ (d) \textit{Penicillium} inhibiting bacterial growth
\\(A) 1-a, 2-b, 3-c, 4-d
\\(B) 1-b, 2-a, 3-d, 4-c
\\(C) 1-c, 2-d, 3-a, 4-b
\\(D) 1-d, 2-c, 3-b, 4-a

\item Population density is measured as:
\\(A) Total number of individuals in an area
\\(B) Number of individuals per unit area or volume
\\(C) Birth rate minus death rate
\\(D) Ratio of juveniles to adults

\end{enumerate}
\end{document}
